% Spoken scripts for the patch-representation section, based on the recording.

\newcommand{\PatchesConditioningScript}{%
{\color{talkblue}\bfseries Slide 13: How much of the dual should we reveal?}\par
{\small 11:00--11:50 (50 seconds)}\par\medskip
First focus on the signed exponential weight. \keyscript{We want a sigma-algebra $\mathcal G_t$ with two properties: the conditional mean should be nonnegative, and it should be explicit enough to analyse.}

If we reveal nothing, we average all the cancellations. The mean is nonnegative: it is $P_t\chi_A$ evaluated at the all-one configuration. But we have gained no way to compute it.

If we reveal the entire dual path, the weight is known, but it still carries its random sign. \keyscript{We need an intermediate amount of information.}%
}

\newcommand{\PatchesGeometryScript}{%
{\color{talkblue}\bfseries Slide 14: The interactions define the patches}\par
{\small 11:50--12:50 (60 seconds)}\par\medskip
Here is our choice. In the graphical construction, births and splits are interactions between sites; deaths are not. \keyscript{Let $\mathcal G_t$ be the sigma-algebra generated by the spacetime positions of the interactions and the sites involved, up to time t.} \keyscript{It does not reveal whether an interaction is a birth or a split, or any deaths.}

The left picture shows the dual realization. On each site line, the incoming and outgoing interactions divide time into intervals. These are the patches shown on the right.

There are four bulk types, according to whether the interaction at each end is incoming or outgoing. Patches reaching the observation time are cut there and become end patches. The thick red line separates bulk patches from end patches.%
}

\newcommand{\PatchesIndependenceScript}{%
{\color{talkblue}\bfseries Slide 15: Conditional independence across patches}\par
{\small 12:50--13:35 (45 seconds)}\par\medskip
The sigma-algebra $\mathcal G_t$ has the properties we wanted. \keyscript{First, it provides enough information for explicit computation: conditional on $\mathcal G_t$, the law of the dual factorizes over patches.} The remaining randomness in different patches is independent.

\keyscript{Each patch has an explicit one-site law $\pP_P^{\mathrm{con}}$, which is why $\mathcal G_t$ satisfies the second condition.}%
}

\newcommand{\PatchesInteriorsScript}{%
{\color{talkblue}\bfseries Slide 16: The four bulk patch types}\par
{\small 13:35--13:45 (10 seconds)}\par\medskip
This figure shows what the process can do inside each patch type. \keyscript{The possibilities are simple enough to compute each patch law explicitly.}%
}

\newcommand{\PatchesRepresentationScript}{%
{\color{talkblue}\bfseries Slide 17: The patch representation}\par
{\small 13:45--15:00 (75 seconds)}\par\medskip
The signed Feynman--Kac weight is itself a product over patches. \keyscript{Conditional independence with respect to $\mathcal G_t$ lets us average each factor separately.} A bulk patch contributes a number $C(P)$. An end patch contributes $C(z,P)$, where z is the initial spin at that site in the original system.

We call the product of these contributions $W_t^\mu$. The initial law $\mu$ acts on the whole product of end contributions, since the initial spins may be correlated.

\keyscript{We can give all these factors explicitly, and all dependence on the initial law remains in the end patches.} This is the patch representation. Next we use patch positivity to control the signs of these contributions.%
}
